% ======================================================================
% BANQUE DE CONTENUS — CV Lucas Vacquié
% ======================================================================
% Chaque macro dessine UN bloc (expérience / engagement / projet) à partir
% d'une coordonnée nord-ouest donnée en argument, et laisse toujours une
% coordonnée de sortie au même nom générique pour pouvoir chaîner les blocs
% choisis facilement dans le CV principal :
%   - Expérience            -> sortie : (exp_last)
%   - Engagement associatif  -> sortie : (asso_last)
%   - Projets réalisés       -> sortie : (proj_last)
%
% Dans le CV principal, tu n'as qu'à commenter/décommenter les lignes
% correspondant aux blocs que tu veux garder pour une candidature donnée.
% Les couleurs / polices (bleu_foncee, gris_foncee, \ptsans, etc.) doivent
% déjà être définies dans le fichier principal AVANT d'inclure ce fichier.
% ======================================================================


% ======================================================================
% SECTION : EXPÉRIENCE   (sortie : (exp_last))
% ======================================================================

% Stage 2ème année ICAM, Assistant Ingénieur, Microtec
\newcommand{\expMicrotec}[1]{%
    \node[anchor=north west, text width=11.0cm, align=justify, font=\ptsansbold\fontsize{10}{8}\selectfont] (exp_sub) at (#1) {
        Assistant Ingénieur (Stage)\\
        Microtec
        \ptsans - Ramonville
    };
    \coordinate (exp_desc_pos) at ($(exp_sub.south west) + (0cm, 0.05cm)$);
    \node[anchor=north west, text width=11.0cm, align=justify, font=\ptsans\fontsize{9}{11}\selectfont, text=gris_foncee] (exp_last) at (exp_desc_pos) {
        \begin{minipage}{\textwidth}
            \begin{itemize}[left=0em, labelsep=0.5em, itemsep=0.1em, align=left, label=\textbullet]
                \item \textbf{Conception d'un outil d'analyse sous Excel (VBA)} pour la qualification par similarité
                \item Implémentation des règles normatives pour valider les composants par héritage
            \end{itemize}
        \end{minipage}
    };
}

%Stage 1er année ICAM, Opérateur de Production, Pliage Métal
\newcommand{\expPliageMetal}[1]{%
    \node[anchor=north west, text width=11.0cm, align=justify,font=\ptsansbold\fontsize{10}{8}\selectfont] (exp_sub) at (#1) {
        Opérateur de Production (Stage Ouvrier)\\
        Pliage Métal
        \ptsans - Montrabé
    };
    \coordinate (exp_desc_pos) at ($(exp_sub.south west) + (0cm, 0.05cm)$);
    \node[anchor=north west, text width=11.0cm, align=justify,font=\ptsans\fontsize{9}{11}\selectfont,text=gris_foncee] (exp_last) at (exp_desc_pos) {
        \begin{minipage}{\textwidth}
            \begin{itemize}[left=0em, labelsep=0.5em, itemsep=0.1em, align=left, label=\textbullet]
                \item  Travail sur machines de précision et respect des processus industriels
                \item  Immersion en atelier de production et respect des normes de sécurité
            \end{itemize}
        \end{minipage}
    };
}

% -- Gabarit pour une future expérience (stage de fin d'études, alternance...) --
% -- À dupliquer / renommer / compléter quand tu as une nouvelle expérience --
\newcommand{\expGabarit}[1]{%
    \node[anchor=north west, text width=11.0cm, align=justify,font=\ptsansbold\fontsize{10}{8}\selectfont] (exp_sub) at (#1) {
        Intitulé du poste (Stage)\\
        Nom Entreprise
        \ptsans - Ville
    };
    \coordinate (exp_desc_pos) at ($(exp_sub.south west) + (0cm, 0.05cm)$);
    \node[anchor=north west, text width=11.0cm, align=justify,font=\ptsans\fontsize{9}{11}\selectfont,text=gris_foncee] (exp_last) at (exp_desc_pos) {
        \begin{minipage}{\textwidth}
            \begin{itemize}[left=0em, labelsep=0.5em, itemsep=0.1em, align=left, label=\textbullet]
                \item  [BROUILLON à compléter]
                \item  [BROUILLON à compléter]
            \end{itemize}
        \end{minipage}
    };
}


% ======================================================================
% SECTION : ENGAGEMENT ASSOCIATIF   (sortie : (asso_last))
% ======================================================================

% Handball, Entraîneur & Arbitre, ASU
\newcommand{\assoHandball}[1]{%
    \node[anchor=north west, text width=11.0cm, align=justify,font=\ptsansbold\fontsize{10}{10}\selectfont] (asso_sub) at (#1) {
        Entraîneur \& Arbitre\\
        ASU Handball
        \ptsans - L'Union
    };
    \coordinate (asso_desc_pos) at ($(asso_sub.south west) + (0cm, 0.05cm)$);
    \node[anchor=north west, text width=11.0cm, align=justify,font=\ptsans\fontsize{9}{11}\selectfont,text=gris_foncee] (asso_last) at (asso_desc_pos) {
        \begin{minipage}{\textwidth}
            \begin{itemize}[left=0em, labelsep=0.5em, itemsep=0.1em, align=left, label=\textbullet]
                \item  Encadrement de jeunes joueurs, organisation de tournois et arbitrage
            \end{itemize}
        \end{minipage}
    };
}

% Projet d'Action Solidaire : Bénévole, Voir et Comprendre
\newcommand{\assoBenevole}[1]{%
    \node[anchor=north west, text width=11.0cm, align=justify,font=\ptsansbold\fontsize{10}{10}\selectfont] (asso_sub) at (#1) {
        Bénévole\\
        Voir et Comprendre
        \ptsans - Toulouse
    };
    \coordinate (asso_desc_pos) at ($(asso_sub.south west) + (0cm, 0.05cm)$);
    \node[anchor=north west, text width=11.0cm, align=justify,font=\ptsans\fontsize{9}{11}\selectfont,text=gris_foncee] (asso_last) at (asso_desc_pos) {
        \begin{minipage}{\textwidth}
            \begin{itemize}[left=0em, labelsep=0.5em, itemsep=0.1em, align=left, label=\textbullet]
                \item  Accompagnement scolaire de collégiens et lycéens (quartier Reynerie)
            \end{itemize}
        \end{minipage}
    };
}

% -- Gabarit pour un futur engagement associatif --
\newcommand{\assoGabarit}[1]{%
    \node[anchor=north west, text width=11.0cm, align=justify,font=\ptsansbold\fontsize{10}{10}\selectfont] (asso_sub) at (#1) {
        Rôle\\
        Nom Association
        \ptsans - Ville
    };
    \coordinate (asso_desc_pos) at ($(asso_sub.south west) + (0cm, 0.05cm)$);
    \node[anchor=north west, text width=11.0cm, align=justify,font=\ptsans\fontsize{9}{11}\selectfont,text=gris_foncee] (asso_last) at (asso_desc_pos) {
        \begin{minipage}{\textwidth}
            \begin{itemize}[left=0em, labelsep=0.5em, itemsep=0.1em, align=left, label=\textbullet]
                \item  [BROUILLON à compléter]
            \end{itemize}
        \end{minipage}
    };
}


% ======================================================================
% SECTION : PROJETS RÉALISÉS   (sortie : (proj_last))
% ======================================================================
% Chaque macro reproduit le style d'origine : un seul node text width=11.0cm
% contenant un itemize à une puce (titre + description).

% Projet Industriel 2ème année ICAM, Système Satellitaire
\newcommand{\projSatellite}[1]{%
    \node[anchor=north west,text width=11.0cm, align=justify,font=\ptsansbold\fontsize{9}{10.5}\selectfont,text=gris_foncee] (proj_last) at (#1) {
        \begin{minipage}{\linewidth}
            \begin{itemize}[left=0em, labelsep=0.5em, itemsep=0.05em, align=left, label=\textbullet]
                \item \href{https://github.com/lucasvcq/satellite-docking-robot}{Système Satellitaire (Architecture Complexe) - Équipe de 25 :}
                \ptsans Conception système d'éjection et prise de vue. Développement de la logique de contrôle et programmation des capteurs (Arduino/ESP32).
            \end{itemize}
        \end{minipage}
    };
}

% Projet Mécatronique 3ème année ICAM, BB8
\newcommand{\projBBEight}[1]{%
    \node[anchor=north west, text width=11.0cm, align=justify, font=\ptsansbold\fontsize{9}{10.5}\selectfont, text=gris_foncee] (proj_last) at (#1) {
        \begin{minipage}{\textwidth}
            \begin{itemize}[left=0em, labelsep=0.5em, itemsep=0.05em, align=left, label=\textbullet]
                \item \href{https://github.com/lucasvcq/BB8-ballbot-truncated-sphere}{Robot Sphérique BB8 (Équipe de 6) :}
                \ptsans Modélisation sous Matlab. Conception électronique et développement de l'asservissement PID et programmation en C++.
            \end{itemize}
        \end{minipage}
    };
}

% Projet 3ème année ICAM, Organisation Nuit ICAM
\newcommand{\projNuitICAM}[1]{%
    \node[anchor=north west, text width=11.0cm, align=justify, font=\ptsansbold\fontsize{9}{10.5}\selectfont, text=gris_foncee] (proj_last) at (#1) {
        \begin{minipage}{\textwidth}
            \begin{itemize}[left=0em, labelsep=0.5em, itemsep=0.05em, align=left, label=\textbullet]
                \item Organisation Nuit ICAM (Gala 1800 pers.) - Pôle Logistique :
                \ptsans Coordination d'une équipe de 80 personnes. Gestion des installations techniques (plan électrique, structures) et supervision opérationnelle.
            \end{itemize}
        \end{minipage}
    };
}

% Experiment entre 2ème et 3ème année ICAM, Mobilité Internationale
\newcommand{\projExperiment}[1]{%
    \node[anchor=north west, text width=11.0cm, align=justify, font=\ptsansbold\fontsize{9}{10.5}\selectfont, text=gris_foncee] (proj_last) at (#1) {
        \begin{minipage}{\textwidth}
            \begin{itemize}[left=0em, labelsep=0.5em, itemsep=0.05em, align=left, label=\textbullet]
                \item "Expériment" - Mobilité Internationale (4 mois, Asie) :
                \ptsans Immersion en autonomie au Cambodge et en Malaisie. 1 mois de bénévolat dans une école et développement de l'adaptabilité interculturelle.
            \end{itemize}
        \end{minipage}
    };
}

% Projet Majeur 4ème année ICAM, Bras Robotique Impression 3D, Thales Alenia Space, Adaxis, RoboDK, UR5e
\newcommand{\projBrasRobotique}[1]{%
    \node[anchor=north west, text width=11.0cm, align=justify, font=\ptsansbold\fontsize{9}{10.5}\selectfont, text=gris_foncee] (proj_last) at (#1) {
        \begin{minipage}{\textwidth}
            \begin{itemize}[left=0em, labelsep=0.5em, itemsep=0.05em, align=left, label=\textbullet]
                \item \href{https://github.com/lucasvcq/robotic-3d-printing-arm}{Cellule d'Impression 3D Robotisée (Collab. Thales Alenia Space)} :
                \ptsans Trajectoire et simulation sous RoboDK / suite Adaxis pour bras Universal Robots. Développement d'un script d'orchestration Python (TCP) et firmware Arduino/RAMPS pour la gestion de la tête d'impression et des sécurités thermiques.
            \end{itemize}
        \end{minipage}
    };
}

% Projet Mineur 4ème année ICAM, Convertisseur Buck Synchrone 48V/12V, Application Automobile
\newcommand{\projBuckConverter}[1]{%
    \node[anchor=north west, text width=11.0cm, align=justify, font=\ptsansbold\fontsize{9}{10.5}\selectfont, text=gris_foncee] (proj_last) at (#1) {
        \begin{minipage}{\textwidth}
            \begin{itemize}[left=0em, labelsep=0.5em, itemsep=0.05em, align=left, label=\textbullet]
                \item Convertisseur Buck Synchrone 48V/12V (Application MHEV) :
                \ptsans Étude et conception d'une fonction d'électronique de puissance pour véhicule hybride léger. Dimensionnement des composants et analyse du schéma de commutation, présentés en soutenance technique.
            \end{itemize}
        \end{minipage}
    };
}

% Projet IA 3ème année ICAM, Comptage de File d'Attente par Vision, Raspberry Pi, Teachable Machine, MobileNet, MQTT, ThingsBoard
\newcommand{\projQueueIA}[1]{%
    \node[anchor=north west, text width=11.0cm, align=justify, font=\ptsansbold\fontsize{9}{10.5}\selectfont, text=gris_foncee] (proj_last) at (#1) {
        \begin{minipage}{\textwidth}
            \begin{itemize}[left=0em, labelsep=0.5em, itemsep=0.05em, align=left, label=\textbullet]
                \item \href{https://github.com/lucasvcq/ai-queue-counting-raspberry-pi}{Système de Comptage de File d'Attente par IA :}
                \ptsans Détection et comptage de personnes par vision (Teachable Machine/MobileNet) sur Raspberry Pi, remontée des données via MQTT vers une plateforme de supervision ThingsBoard.
            \end{itemize}
        \end{minipage}
    };
}

% Projet Jumeaux Numériques 4ème année ICAM, Robot Soccer Kit, Python, Simulation, Modélisation
\newcommand{\projSoccerKit}[1]{%
    \node[anchor=north west, text width=11.0cm, align=justify, font=\ptsansbold\fontsize{9}{10.5}\selectfont, text=gris_foncee] (proj_last) at (#1) {
        \begin{minipage}{\textwidth}
            \begin{itemize}[left=0em, labelsep=0.5em, itemsep=0.05em, align=left, label=\textbullet]
                \item \href{https://github.com/lucasvcq/robot-soccer-digital-twin}{Jumeau Numérique de Robot Soccer (Robot Soccer Kit, Python) :}
                \ptsans Modélisation et simulation en Python du comportement de robots de compétition, pour tester des stratégies avant déploiement réel.
            \end{itemize}
        \end{minipage}
    };
}

% Projet 4ème année ICAM, Application Mobile de Vie Associative, React Native, Expo, TypeScript, Firebase, Gemini API
\newcommand{\projAppMobileAsso}[1]{%
    \node[anchor=north west, text width=11.0cm, align=justify, font=\ptsansbold\fontsize{9}{10.5}\selectfont, text=gris_foncee] (proj_last) at (#1) {
        \begin{minipage}{\textwidth}
            \begin{itemize}[left=0em, labelsep=0.5em, itemsep=0.05em, align=left, label=\textbullet]
                \item \href{https://github.com/lucasvcq/student-sport-app}{Application Web/Mobile de Vie Associative ICAM (Équipe de 3) :}
                \ptsans Développement d'un POC full-stack (\textbf{React Native / Expo}, \textbf{TypeScript}, \textbf{Firebase}). Intégration d'un assistant IA conversationnel (\textbf{Gemini API} avec Function Calling), gestion des rôles/RSVP et génération automatique de documents PDF.
            \end{itemize}
        \end{minipage}
    };
}

% -- Gabarit pour un futur projet --
\newcommand{\projGabarit}[1]{%
    \node[anchor=north west, text width=11.0cm, align=justify, font=\ptsansbold\fontsize{9}{10.5}\selectfont, text=gris_foncee] (proj_last) at (#1) {
        \begin{minipage}{\textwidth}
            \begin{itemize}[left=0em, labelsep=0.5em, itemsep=0.05em, align=left, label=\textbullet]
                \item Titre du Projet (Contexte, taille d'équipe) :
                \ptsans [BROUILLON à compléter]
            \end{itemize}
        \end{minipage}
    };
}
